\title{2023 年高考数学 天津卷}
\date{}
\maketitle

\section*{一、选择题}
在每小题给出的四个选项中，只有一项是符合题目要求的

\begin{question}[points=5]
已知集合 $U = \{1, 2, 3, 4, 5\}$, $A = \{1, 3\}$, $B = \{1, 2, 4\}$, 则 $\complement_U B \cup A =$ ( )
\begin{choices}
  \item $\{1,3,5\}$
  \item $\{1,3\}$
  \item $\{1,2,4\}$
  \item $\{1,2,4,5\}$
\end{choices}
\begin{solution}
\textbf{答案：}A

\textbf{分析：}对集合 $B$ 求补集，应用集合的并运算求结果.

\textbf{详解：}由 $\complement_U B = \{3,5\}$, 而 $A = \{1,3\}$, 所以 $\complement_U B \cup A = \{1,3,5\}$.
\end{solution}
\end{question}

\begin{question}[points=5]
“$a^2 = b^2$”是“$a^2 + b^2 = 2ab$”的 ( )
\begin{choices}
  \item 充分不必要条件
  \item 必要不充分条件
  \item 充分必要条件
  \item 既不充分又不必要条件
\end{choices}
\begin{solution}
\textbf{答案：}B

\textbf{分析：}根据充分、必要性定义判断条件的推出关系.

\textbf{详解：}由 $a^2 = b^2$, 则 $a = \pm b$, 当 $a = -b \neq 0$ 时 $a^2 + b^2 = 2ab$ 不成立, 充分性不成立. 由 $a^2 + b^2 = 2ab$, 则 $(a-b)^2 = 0$, 即 $a = b$, 显然 $a^2 = b^2$ 成立, 必要性成立. 所以 $a^2 = b^2$ 是 $a^2 + b^2 = 2ab$ 的必要不充分条件.
\end{solution}
\end{question}

\begin{question}[points=5]
若 $a = 1.01^{0.5}$, $b = 1.01^{0.6}$, $c = 0.6^{0.5}$, 则 $a, b, c$ 的大小关系为 ( )
\begin{choices}
  \item $c > a > b$
  \item $c > b > a$
  \item $a > b > c$
  \item $b > a > c$
\end{choices}
\begin{solution}
\textbf{答案：}D

\textbf{分析：}根据对应幂、指数函数的单调性判断大小关系.

\textbf{详解：}由 $y = 1.01^x$ 在 $\mathbb{R}$ 上递增, 则 $a = 1.01^{0.5} < b = 1.01^{0.6}$. 由 $y = x^{0.5}$ 在 $[0, +\infty)$ 上递增, 则 $a = 1.01^{0.5} > c = 0.6^{0.5}$. 所以 $b > a > c$.
\end{solution}
\end{question}

\begin{question}[points=5]
函数 $f(x)$ 的图象如下图所示, 则 $f(x)$ 的解析式可能为 ( )

\begin{center}
  \includegraphics[width=.38\linewidth]{question-4-1.jpg}
\end{center}
\begin{choices}
  \item $\dfrac{5(e^x - e^{-x})}{x^2 + 2}$
  \item $\dfrac{5\sin x}{x^2 + 1}$
  \item $\dfrac{5(e^x + e^{-x})}{x^2 + 2}$
  \item $\dfrac{5\cos x}{x^2 + 1}$
\end{choices}
\begin{solution}
\textbf{答案：}D

\textbf{分析：}由图知函数为偶函数, 应用排除法.

\textbf{详解：}由图知: 函数图象关于 $y$ 轴对称, 其为偶函数, 且 $f(-2) = f(2) < 0$. 由 $\dfrac{5\sin(-x)}{(-x)^2 + 1} = -\dfrac{5\sin x}{x^2 + 1}$ 且定义域为 $\mathbb{R}$, 即 B 中函数为奇函数, 排除. 当 $x > 0$ 时 $\dfrac{5(e^x - e^{-x})}{x^2 + 2} > 0$, $\dfrac{5(e^x + e^{-x})}{x^2 + 2} > 0$, 即 A、C 中 $(0, +\infty)$ 上函数值为正, 排除. 故选 D.
\end{solution}
\end{question}

\begin{question}[points=5]
已知函数 $f(x)$ 的一条对称轴为直线 $x = 2$, 一个周期为 $4$, 则 $f(x)$ 的解析式可能为 ( )
\begin{choices}
  \item $\sin\left(\dfrac{\pi}{2}x\right)$
  \item $\cos\left(\dfrac{\pi}{2}x\right)$
  \item $\sin\left(\dfrac{\pi}{4}x\right)$
  \item $\cos\left(\dfrac{\pi}{4}x\right)$
\end{choices}
\begin{solution}
\textbf{答案：}B

\textbf{分析：}由题意分别考查函数的最小正周期和函数在 $x=2$ 处的函数值.

\textbf{详解：}A 选项 $T = \dfrac{2\pi}{\pi/2} = 4$, B 选项 $T = \dfrac{2\pi}{\pi/2} = 4$, C 选项 $T = \dfrac{2\pi}{\pi/4} = 8$, D 选项 $T = \dfrac{2\pi}{\pi/4} = 8$, 排除 CD. 对于 A, 当 $x=2$ 时 $\sin\left(\dfrac{\pi}{2}\cdot 2\right)=0$, 故 $(2,0)$ 是一个对称中心, 排除 A. 对于 B, 当 $x=2$ 时 $\cos\left(\dfrac{\pi}{2}\cdot 2\right) = -1$, 故 $x=2$ 是一条对称轴. 故选 B.
\end{solution}
\end{question}

\begin{question}[points=5]
已知 $\{a_n\}$ 为等比数列, $S_n$ 为数列 $\{a_n\}$ 的前 $n$ 项和, $a_{n+1} = 2S_n + 2$, 则 $a_4$ 的值为 ( )
\begin{choices}
  \item 3
  \item 18
  \item 54
  \item 152
\end{choices}
\begin{solution}
\textbf{答案：}C

\textbf{分析：}对递推关系式进行赋值, 得到关于首项、公比的方程组.

\textbf{详解：}当 $n=1$ 时, $a_2 = 2a_1 + 2$, 即 $a_1 q = 2a_1 + 2$ ①. 当 $n=2$ 时, $a_3 = 2(a_1 + a_2) + 2$, 即 $a_1 q^2 = 2(a_1 + a_1 q) + 2$ ②. 联立①②得 $a_1 = 2$, $q = 3$, 则 $a_4 = a_1 q^3 = 54$.
\end{solution}
\end{question}

\begin{question}[points=5]
调查某种群花萼长度和花瓣长度, 所得数据如图所示, 其中相关系数 $r = 0.8245$, 下列说法正确的是 ( )

\begin{center}
  \includegraphics[width=.38\linewidth]{question-7-1.jpg}
\end{center}
\begin{choices}
  \item 花瓣长度和花萼长度没有相关性
  \item 花瓣长度和花萼长度呈现负相关
  \item 花瓣长度和花萼长度呈现正相关
  \item 若从样本中抽取一部分, 则这部分的相关系数一定是 $0.8245$
\end{choices}
\begin{solution}
\textbf{答案：}C

\textbf{分析：}根据散点图的特点可分析出相关性, 根据相关系数的定义判断 D.

\textbf{详解：}根据散点的集中程度可知, 花瓣长度和花萼长度有相关性, A 错误. 散点的分布是从左下到右上, 呈现正相关性, B 错误, C 正确. 由于 $r = 0.8245$ 是全部数据的相关系数, 取出来一部分数据, 相关系数可能变化, D 错误.
\end{solution}
\end{question}

\begin{question}[points=5]
在三棱锥 $P-ABC$ 中, 线段 $PC$ 上的点 $M$ 满足 $PM = \dfrac{1}{3}PC$, 线段 $PB$ 上的点 $N$ 满足 $PN = \dfrac{2}{3}PB$, 则三棱锥 $P-AMN$ 和三棱锥 $P-ABC$ 的体积之比为 ( )
\begin{choices}
  \item $\dfrac{1}{9}$
  \item $\dfrac{2}{9}$
  \item $\dfrac{1}{3}$
  \item $\dfrac{4}{9}$
\end{choices}
\begin{solution}
\textbf{答案：}B

\textbf{分析：}分别过 M, C 作 $MM' \perp PA$, $CC' \perp PA$, 过 B 作 $BB' \perp$ 平面 $PAC$, 连接 $PB'$, 过 N 作 $NN' \perp PB'$. 利用相似比和体积公式.

\textbf{详解：}如图, 分别过 $M, C$ 作 $MM' \perp PA$, $CC' \perp PA$, 垂足分别为 $M', C'$. 过 $B$ 作 $BB' \perp$ 平面 $PAC$, 垂足为 $B'$, 连接 $PB'$, 过 $N$ 作 $NN' \perp PB'$, 垂足为 $N'$. 因为 $BB' \perp$ 平面 $PAC$, $BB' \subset$ 平面 $PBB'$, 所以平面 $PBB' \perp$ 平面 $PAC$. 又因为平面 $PBB' \cap$ 平面 $PAC = PB'$, $NN' \perp PB'$, $NN' \subset$ 平面 $PBB'$, 所以 $NN' \perp$ 平面 $PAC$, 且 $BB' \parallel NN'$. 在 $\triangle PCC'$ 中, 因为 $MM' \perp PA$, $CC' \perp PA$, 所以 $MM' \parallel CC'$, 所以 $\dfrac{PM}{PC} = \dfrac{MM'}{CC'} = \dfrac{1}{3}$. 在 $\triangle PBB'$ 中, 因为 $BB' \parallel NN'$, 所以 $\dfrac{PN}{PB} = \dfrac{NN'}{BB'} = \dfrac{2}{3}$. 所以 $\dfrac{V_{P-AMN}}{V_{P-ABC}} = \dfrac{V_{N-PAM}}{V_{B-PAC}} = \dfrac{\frac{1}{3} S_{\triangle PAM} \cdot NN'}{\frac{1}{3} S_{\triangle PAC} \cdot BB'} = \dfrac{\frac{1}{2} PA \cdot MM' \cdot NN'}{\frac{1}{2} PA \cdot CC' \cdot BB'} = \dfrac{1}{3} \cdot \frac{2}{3} = \dfrac{2}{9}$.
\end{solution}
\end{question}

\begin{question}[points=5]
双曲线 $\dfrac{x^2}{a^2} - \dfrac{y^2}{b^2} = 1$ ($a > 0$, $b > 0$) 的左、右焦点分别为 $F_1$, $F_2$. 过 $F_2$ 作其中一条渐近线的垂线, 垂足为 $P$. 已知 $PF_2 = 2$, 直线 $PF_1$ 的斜率为 $\dfrac{\sqrt{2}}{4}$, 则双曲线的方程为 ( )
\begin{choices}
  \item $\dfrac{x^2}{8} - \dfrac{y^2}{4} = 1$
  \item $\dfrac{x^2}{4} - \dfrac{y^2}{8} = 1$
  \item $\dfrac{x^2}{4} - \dfrac{y^2}{2} = 1$
  \item $\dfrac{x^2}{2} - \dfrac{y^2}{4} = 1$
\end{choices}
\begin{solution}
\textbf{答案：}D

\textbf{分析：}由点到直线的距离公式求出 $b$, 设 $\angle POF_2 = \theta$, 由 $\tan \theta = \frac{b}{a}$ 得到 $OP = a$, $OF_2 = c$. 再由三角形面积公式和斜率公式.

\textbf{详解：}因为 $F_2(c,0)$, 不妨设渐近线方程为 $y = \frac{b}{a}x$, 即 $bx - ay = 0$, 所以 $PF_2 = \frac{bc}{\sqrt{a^2 + b^2}} = \frac{bc}{c} = b = 2$. 设 $\angle POF_2 = \theta$, 则 $\tan \theta = \frac{PF_2}{OP} = \frac{b}{OP} = \frac{b}{a}$, 所以 $OP = a$, $OF_2 = c$. 因为 $\frac{1}{2}ab = \frac{1}{2}c \cdot y_P$, 所以 $y_P = \frac{ab}{c}$, 又 $\tan \theta = \frac{y_P}{x_P} = \frac{b}{a}$, 所以 $x_P = \frac{a^2}{c}$. 所以 $P\left(\frac{a^2}{c}, \frac{ab}{c}\right)$. 因为 $F_1(-c,0)$, 所以 $k_{PF_1} = \frac{\frac{ab}{c}}{\frac{a^2}{c} + c} = \frac{ab}{a^2 + c^2} = \frac{a \cdot 2}{a^2 + a^2 + 4} = \frac{2a}{2a^2 + 4} = \frac{a}{a^2 + 2} = \frac{\sqrt{2}}{4}$, 所以 $4a = \sqrt{2}(a^2 + 2)$, 解得 $a = \sqrt{2}$, 所以双曲线方程为 $\frac{x^2}{2} - \frac{y^2}{4} = 1$.
\end{solution}
\end{question}

\section*{二、填空题}
本大题共 6 小题，每小题 5 分，共 30 分．试题中包含两个空的，答对 1 个的给 3 分，全部答对的给 5 分．

\begin{question}[points=5]
已知 $i$ 是虚数单位，化简 $\dfrac{5 + 14i}{2 + 3i}$ 的结果为．
\fillin[{$4 + i$}]
\begin{solution}
\textbf{答案：}$4 + i$

\textbf{分析：}分子分母同时乘以 $2 - 3i$.

\textbf{详解：}$\dfrac{5 + 14i}{2 + 3i} = \dfrac{(5+14i)(2-3i)}{(2+3i)(2-3i)} = \dfrac{52 + 13i}{13} = 4 + i$.
\end{solution}
\end{question}

\begin{question}[points=5]
在 $\left(2x^3 - \dfrac{1}{x}\right)^6$ 的展开式中，$x^2$ 项的系数为．
\fillin[{$60$}]
\begin{solution}
\textbf{答案：}$60$

\textbf{分析：}由二项式展开式的通项公式写出通项，令 $x$ 的指数为 $2$.

\textbf{详解：}通项公式 $T_{k+1} = C_6^k (2x^3)^{6-k} \left(-\dfrac{1}{x}\right)^k = (-1)^k 2^{6-k} C_6^k x^{18-4k}$，令 $18-4k=2$，得 $k=4$，则 $x^2$ 项的系数为 $(-1)^4 2^{2} C_6^4 = 4 \times 15 = 60$.
\end{solution}
\end{question}

\begin{question}[points=5]
过原点的一条直线与圆 $C: (x+2)^2 + y^2 = 3$ 相切，交曲线 $y^2 = 2px$ ($p > 0$) 于点 $P$，若 $OP = 8$，则 $p$ 的值为．
\fillin[{$6$}]
\begin{solution}
\textbf{答案：}$6$

\textbf{分析：}根据圆和曲线关于 $x$ 轴对称，设切线方程为 $y = kx$ ($k>0$)，由直线与圆相切得 $k$，再联立抛物线得 $P$ 坐标，由 $OP=8$ 求 $p$.

\textbf{详解：}易知圆 $(x+2)^2 + y^2 = 3$ 和曲线 $y^2 = 2px$ 关于 $x$ 轴对称，不妨设切线方程为 $y = kx$，$k > 0$，则 $\dfrac{|2k|}{\sqrt{1+k^2}} = \sqrt{3}$，解得 $k = \sqrt{3}$. 由 $\begin{cases} y = \sqrt{3}x \\ y^2 = 2px \end{cases}$ 解得 $\begin{cases} x = 0 \\ y = 0 \end{cases}$ 或 $\begin{cases} x = \dfrac{2p}{3} \\ y = \dfrac{2\sqrt{3}p}{3} \end{cases}$. 所以 $OP = \sqrt{\left(\dfrac{2p}{3}\right)^2 + \left(\dfrac{2\sqrt{3}p}{3}\right)^2} = \dfrac{4p}{3} = 8$，解得 $p=6$.
\end{solution}
\end{question}

\begin{question}[points=5]
甲乙丙三个盒子中装有一定数量的黑球和白球，其总数之比为 $5:4:6$．这三个盒子中黑球占总数的比例分别为 $40\%, 25\%, 50\%$．现从三个盒子中各取一个球，取到的三个球都是黑球的概率为；将三个盒子混合后任取一个球，是白球的概率为．
\fillin[{$0.05$ ; $\dfrac{3}{5}$}]
\begin{solution}
\textbf{答案：}$0.05$ ; $\dfrac{3}{5}$

\textbf{分析：}先求出各盒中黑球、白球数量，再用概率乘法公式和古典概型公式.

\textbf{详解：}设甲乙丙三个盒子中的球数分别为 $5n,4n,6n$，总数 $15n$. 甲盒黑球 $40\% \times 5n = 2n$，白球 $3n$；乙盒黑球 $25\% \times 4n = n$，白球 $3n$；丙盒黑球 $50\% \times 6n = 3n$，白球 $3n$. 从三个盒子中各取一个球都是黑球的概率 $P(A) = 0.4 \times 0.25 \times 0.5 = 0.05$. 混合后黑球总数 $2n + n + 3n = 6n$，白球总数 $9n$，所以 $P(B) = \dfrac{9n}{15n} = \dfrac{3}{5}$.
\end{solution}
\end{question}

\begin{question}[points=5]
在 $\triangle ABC$ 中，$\angle A = 60^\circ$，$BC = 1$，点 $D$ 为 $AB$ 的中点，点 $E$ 为 $CD$ 的中点，若设 $\overrightarrow{AB} = \vec{a}$，$\overrightarrow{AC} = \vec{b}$，则 $\overrightarrow{AE}$ 可用 $\vec{a}, \vec{b}$ 表示为；若 $\overrightarrow{BF} = \dfrac{1}{3} \overrightarrow{BC}$，则 $\overrightarrow{AE} \cdot \overrightarrow{AF}$ 的最大值为．
\fillin[{$\dfrac{1}{4}\vec{a} + \dfrac{1}{2}\vec{b}$ ; $\dfrac{13}{24}$}]
\begin{solution}
\textbf{答案：}$\dfrac{1}{4}\vec{a} + \dfrac{1}{2}\vec{b}$ ; $\dfrac{13}{24}$

\textbf{分析：}利用向量线性运算和数量积，结合余弦定理和基本不等式.

\textbf{详解：}因为 $E$ 为 $CD$ 的中点，则 $\overrightarrow{ED} + \overrightarrow{EC} = \vec{0}$，由 $\begin{cases} \overrightarrow{AE} + \overrightarrow{ED} = \overrightarrow{AD} \\ \overrightarrow{AE} + \overrightarrow{EC} = \overrightarrow{AC} \end{cases}$ 两式相加得 $2\overrightarrow{AE} = \overrightarrow{AD} + \overrightarrow{AC} = \dfrac{1}{2}\overrightarrow{AB} + \overrightarrow{AC} = \dfrac{1}{2}\vec{a} + \vec{b}$，所以 $\overrightarrow{AE} = \dfrac{1}{4}\vec{a} + \dfrac{1}{2}\vec{b}$. 因为 $\overrightarrow{BF} = \dfrac{1}{3}\overrightarrow{BC}$，则 $2\overrightarrow{FB} + \overrightarrow{FC} = \vec{0}$，由 $\begin{cases} \overrightarrow{AF} + \overrightarrow{FC} = \overrightarrow{AC} \\ \overrightarrow{AF} + \overrightarrow{FB} = \overrightarrow{AB} \end{cases}$ 两式相加得 $3\overrightarrow{AF} = 2\overrightarrow{AB} + \overrightarrow{AC} = 2\vec{a} + \vec{b}$，所以 $\overrightarrow{AF} = \dfrac{2}{3}\vec{a} + \dfrac{1}{3}\vec{b}$. 所以 $\overrightarrow{AE} \cdot \overrightarrow{AF} = \left(\dfrac{1}{4}\vec{a} + \dfrac{1}{2}\vec{b}\right) \cdot \left(\dfrac{2}{3}\vec{a} + \dfrac{1}{3}\vec{b}\right) = \dfrac{1}{12}(2|\vec{a}|^2 + 5\vec{a}\cdot\vec{b} + 2|\vec{b}|^2)$. 设 $AB = x$, $AC = y$, 则 $\vec{a}\cdot\vec{b} = xy\cos60^\circ = \dfrac{xy}{2}$, 且 $BC^2 = x^2 + y^2 - xy = 1$, 所以 $2x^2 + 5\cdot\frac{xy}{2} + 2y^2 = 2(x^2 + y^2) + \frac{5}{2}xy = 2(1+xy) + \frac{5}{2}xy = 2 + \frac{9}{2}xy$. 由 $x^2 + y^2 \ge 2xy$ 得 $1 = x^2 + y^2 - xy \ge xy$, 即 $xy \le 1$, 当 $x = y = 1$ 时取等. 所以 $\overrightarrow{AE} \cdot \overrightarrow{AF} = \dfrac{1}{12}\left(2 + \frac{9}{2}xy\right) \le \dfrac{1}{12}\left(2 + \frac{9}{2}\right) = \dfrac{13}{24}$.
\end{solution}
\end{question}

\begin{question}[points=5]
若函数 $f(x) = |ax^2 - 2x| - |x^2 - ax + 1|$ 有且仅有两个零点，则 $a$ 的取值范围为．
\fillin[{$(-\infty, 0) \cup (0, 1) \cup (1, +\infty)$}]
\begin{solution}
\textbf{答案：}$(-\infty, 0) \cup (0, 1) \cup (1, +\infty)$

\textbf{分析：}分类讨论去掉绝对值，求出零点，再根据根存在的条件判断.

\textbf{详解：}当 $x^2 - ax + 1 \ge 0$ 时, $f(x)=0 \Leftrightarrow (a-1)x^2 + (a-2)x - 1 = 0$, 即 $[(a-1)x - 1](x+1)=0$. 若 $a=1$, 则 $x=-1$; 若 $a\neq1$, 则 $x = \frac{1}{a-1}$ 或 $x=-1$. 条件要求 $x^2 - ax + 1 \ge 0$ 成立, 对于 $x=-1$, 需 $1+a+1\ge0$ 即 $a\ge -2$; 对于 $x = \frac{1}{a-1}$, 需 $\frac{1}{(a-1)^2} - \frac{a}{a-1}+1 \ge 0$ 即 $a \le 2$ 且 $a\neq1$. 当 $x^2 - ax + 1 < 0$ 时, $f(x)=0 \Leftrightarrow (a+1)x^2 - (a+2)x + 1 = 0$, 即 $[(a+1)x - 1](x-1)=0$. 若 $a=-1$, 则 $x=1$ 不满足; 若 $a\neq -1$, 则 $x=1$ 或 $x = \frac{1}{a+1}$. 条件要求 $1 - a + 1 < 0$ 即 $a>2$ 对于 $x=1$, 对于 $x = \frac{1}{a+1}$ 需 $\frac{1}{(a+1)^2} - \frac{a}{a+1} + 1 < 0$ 即 $a<-2$. 综合所有情况, 当 $a<-2$, 零点为 $\frac{1}{a+1}$ 和 $\frac{1}{a-1}$; 当 $-2\le a<0$, 零点为 $\frac{1}{a-1}$ 和 $-1$; 当 $a=0$, 只有一个零点 $-1$; 当 $0<a<1$, 零点为 $\frac{1}{a-1}$ 和 $-1$; 当 $a=1$, 只有一个零点 $-1$; 当 $1<a\le2$, 零点为 $\frac{1}{a-1}$ 和 $-1$; 当 $a>2$, 零点为 $1$ 和 $-1$. 所以有两个零点时, $a\neq0$ 且 $a\neq1$, 即 $a \in (-\infty,0) \cup (0,1) \cup (1,+\infty)$.
\end{solution}
\end{question}

\section*{三、解答题}
本大题共 5 小题，共 75 分，解答应写出文字说明，证明过程或演算步骤．

\begin{question}[points=14]
在 $\triangle ABC$ 中，角 $A, B, C$ 所对的边分别是 $a, b, c$．已知 $a = \sqrt{39}$, $b = 2$, $\angle A = 120^\circ$．

(1) 求 $\sin B$ 的值；

(2) 求 $c$ 的值；

(3) 求 $\sin(B - C)$．
\begin{solution}
\textbf{答案：}(1) $\dfrac{\sqrt{13}}{13}$ ; (2) $5$ ; (3) $-\dfrac{7\sqrt{3}}{26}$

\textbf{详解：}(1) 由正弦定理 $\dfrac{a}{\sin A} = \dfrac{b}{\sin B}$, 得 $\dfrac{\sqrt{39}}{\sin120^\circ} = \dfrac{2}{\sin B}$, 解得 $\sin B = \dfrac{\sqrt{13}}{13}$.
(2) 由余弦定理 $a^2 = b^2 + c^2 - 2bc\cos A$, 即 $39 = 4 + c^2 - 2\cdot 2 \cdot c \cdot \left(-\frac{1}{2}\right) = 4 + c^2 + 2c$, 整理得 $c^2 + 2c - 35 = 0$, 解得 $c = 5$ (负值舍去).
(3) 由正弦定理 $\dfrac{a}{\sin A} = \dfrac{c}{\sin C}$ 得 $\dfrac{\sqrt{39}}{\sin120^\circ} = \dfrac{5}{\sin C}$, 解得 $\sin C = \dfrac{5\sqrt{13}}{26}$. 因为 $A = 120^\circ$, 所以 $B, C$ 均为锐角, $\cos B = \sqrt{1-\left(\frac{\sqrt{13}}{13}\right)^2} = \dfrac{2\sqrt{39}}{13}$, $\cos C = \sqrt{1-\left(\frac{5\sqrt{13}}{26}\right)^2} = \dfrac{\sqrt{39}}{26}$. 所以 $\sin(B-C) = \sin B \cos C - \cos B \sin C = \dfrac{\sqrt{13}}{13} \cdot \dfrac{\sqrt{39}}{26} - \dfrac{2\sqrt{39}}{13} \cdot \dfrac{5\sqrt{13}}{26} = -\dfrac{7\sqrt{3}}{26}$.
\end{solution}
\end{question}

\begin{question}[points=15]
三棱台 $ABC-A_1B_1C_1$ 中，若 $A_1A \perp$ 面 $ABC$, $AB \perp AC$, $AB = AC = AA_1 = 2$, $A_1C_1 = 1$, $M, N$ 分别是 $BC, BA$ 中点.

(1) 求证：$A_1N \parallel$ 平面 $C_1MA$;

(2) 求平面 $C_1MA$ 与平面 $ACC_1A_1$ 所成夹角的余弦值;

(3) 求点 $C$ 到平面 $C_1MA$ 的距离.

\begin{center}
  \includegraphics[width=.38\linewidth]{question-17-1.jpg}
\end{center}
\begin{solution}
\textbf{答案：}(1) 证明见解析; (2) $\dfrac{2}{3}$; (3) $\dfrac{4}{3}$

\textbf{详解：}(1) 连接 $MN$, $C_1A$. 由 $M,N$ 分别是 $BC,BA$ 中点, 得 $MN \parallel AC$ 且 $MN = \frac{AC}{2}=1$. 由棱台性质, $A_1C_1 \parallel AC$, 所以 $MN \parallel A_1C_1$, 且 $MN = A_1C_1 = 1$, 故四边形 $MNA_1C_1$ 是平行四边形, 所以 $A_1N \parallel MC_1$. 又 $A_1N \not\subset$ 平面 $C_1MA$, $MC_1 \subset$ 平面 $C_1MA$, 所以 $A_1N \parallel$ 平面 $C_1MA$.
(2) 过 $M$ 作 $ME \perp AC$ 于 $E$, 过 $E$ 作 $EF \perp AC_1$ 于 $F$, 连接 $MF$. 由 $A_1A \perp$ 面 $ABC$, $ME \subset$ 面 $ABC$, 得 $AA_1 \perp ME$, 又 $ME \perp AC$, $AC \cap AA_1 = A$, 所以 $ME \perp$ 平面 $ACC_1A_1$. 故 $ME \perp AC_1$. 又 $EF \perp AC_1$, 所以 $AC_1 \perp$ 平面 $MEF$, 则 $AC_1 \perp MF$. 故 $\angle MFE$ 为平面 $C_1MA$ 与平面 $ACC_1A_1$ 所成角. 计算得 $ME = \frac{AB}{2}=1$, $\cos\angle CAC_1 = \frac{1}{\sqrt{5}}$, $\sin\angle CAC_1 = \frac{2}{\sqrt{5}}$, $EF = 1 \cdot \sin\angle CAC_1 = \frac{2}{\sqrt{5}}$, $MF = \sqrt{ME^2 + EF^2} = \sqrt{1 + \frac{4}{5}} = \frac{3}{\sqrt{5}}$, 所以 $\cos\angle MFE = \frac{EF}{MF} = \frac{2}{3}$.
(3) [方法一] 过 $C_1$ 作 $C_1P \perp AC$ 于 $P$, 作 $C_1Q \perp AM$ 于 $Q$, 连接 $PQ, PM$, 过 $P$ 作 $PR \perp C_1Q$ 于 $R$. 计算得 $C_1A = C_1C = \sqrt{5}$, $C_1M = \sqrt{5}$, $C_1Q = \frac{3\sqrt{2}}{2}$, $PR = \frac{PC_1 \cdot PQ}{QC_1} = \frac{2 \cdot \frac{\sqrt{2}}{2}}{\frac{3\sqrt{2}}{2}} = \frac{2}{3}$. 因为 $CA = 2PA$, 所以点 $C$ 到平面 $C_1MA$ 的距离是点 $P$ 到该平面距离的两倍, 即 $\frac{4}{3}$.
[方法二] 设点 $C$ 到平面 $C_1MA$ 的距离为 $h$. $V_{C_1-AMC} = \frac{1}{3} \cdot C_1P \cdot S_{\triangle AMC} = \frac{1}{3} \cdot 2 \cdot \frac{1}{2} \cdot 2 \cdot 2 = \frac{4}{3}$. 又 $V_{C-C_1MA} = \frac{1}{3} \cdot h \cdot S_{\triangle AMC_1} = \frac{1}{3} \cdot h \cdot \frac{1}{2} \cdot 2 \cdot \frac{3\sqrt{2}}{2} = \frac{h}{2}$. 由 $V_{C_1-AMC} = V_{C-C_1MA}$ 得 $\frac{h}{2} = \frac{4}{3}$, 所以 $h = \frac{4}{3}$.
\end{solution}
\end{question}

\begin{question}[points=15]
设椭圆 $\dfrac{x^2}{a^2} + \dfrac{y^2}{b^2} = 1$ ($a > b > 0$) 的左右顶点分别为 $A_1$, $A_2$, 右焦点为 $F$, 已知 $A_1F = 3$, $A_2F = 1$.

(1) 求椭圆方程及其离心率;

(2) 已知点 $P$ 是椭圆上一动点 (不与端点重合), 直线 $A_2P$ 交 $y$ 轴于点 $Q$, 若三角形 $A_1PQ$ 的面积是三角形 $A_2FP$ 面积的二倍, 求直线 $A_2P$ 的方程.
\begin{solution}
\textbf{答案：}(1) $\dfrac{x^2}{4} + \dfrac{y^2}{3} = 1$, $e = \dfrac{1}{2}$; (2) $y = \pm \dfrac{\sqrt{6}}{2}(x-2)$

\textbf{详解：}(1) 由题意得 $\begin{cases} a + c = 3 \\ a - c = 1 \end{cases}$, 解得 $a=2$, $c=1$, 所以 $b = \sqrt{2^2 - 1^2} = \sqrt{3}$, 椭圆方程为 $\dfrac{x^2}{4} + \dfrac{y^2}{3} = 1$, 离心率 $e = \dfrac{c}{a} = \dfrac{1}{2}$.
(2) 由题意 $A_2(2,0)$, 设直线 $A_2P$ 的方程为 $y = k(x-2)$, $k$ 存在. 联立 $\begin{cases} \dfrac{x^2}{4} + \dfrac{y^2}{3} = 1 \\ y = k(x-2) \end{cases}$, 消去 $y$ 得 $(3+4k^2)x^2 - 16k^2 x + 16k^2 - 12 = 0$. 由韦达定理 $x_{A_2} x_P = \dfrac{16k^2 - 12}{3+4k^2}$, 所以 $x_P = \dfrac{8k^2 - 6}{3+4k^2}$, 代入得 $P\left(\dfrac{8k^2 - 6}{3+4k^2}, \dfrac{-12k}{3+4k^2}\right)$, $Q(0, -2k)$. 则 $S_{\triangle A_2QA_1} = \frac{1}{2} \cdot 4 \cdot |y_Q|$, $S_{\triangle A_2FP} = \frac{1}{2} \cdot 1 \cdot |y_P|$, $S_{\triangle A_1A_2P} = \frac{1}{2} \cdot 4 \cdot |y_P|$. 由 $S_{\triangle A_1PQ} = S_{\triangle A_2QA_1} - S_{\triangle A_1A_2P} = 2S_{\triangle A_2FP}$, 得 $2|y_Q| = 3|y_P|$, 即 $2|-2k| = 3\left|\dfrac{-12k}{3+4k^2}\right|$, 解得 $k = \pm \dfrac{\sqrt{6}}{2}$. 所以直线 $A_2P$ 的方程为 $y = \pm \dfrac{\sqrt{6}}{2}(x-2)$.
\end{solution}
\end{question}

\begin{question}[points=15]
已知 $\{a_n\}$ 是等差数列, $a_2 + a_5 = 16$, $a_5 - a_3 = 4$.

(1) 求 $\{a_n\}$ 的通项公式和 $\displaystyle\sum_{i=2^{n-1}}^{2^n-1} a_i$;

(2) 已知 $\{b_n\}$ 为等比数列, 对于任意 $k \in \mathbb{N}^*$, 若 $2^{k-1} \le n \le 2^k - 1$, 则 $b_k < a_n < b_{k+1}$,

   (Ⅰ) 当 $k \ge 2$ 时, 求证: $2^k - 1 < b_k < 2^k + 1$;

   (Ⅱ) 求 $\{b_n\}$ 的通项公式及其前 $n$ 项和.
\begin{solution}
\textbf{答案：}(1) $a_n = 2n + 1$, $\displaystyle\sum_{i=2^{n-1}}^{2^n-1} a_i = 3 \cdot 2^{2n-1}$; (2) (Ⅰ) 证明见解析; (Ⅱ) $b_n = 2^n$, 前 $n$ 项和为 $2^{n+1} - 2$

\textbf{详解：}(1) 由 $\begin{cases} a_2 + a_5 = 2a_1 + 5d = 16 \\ a_5 - a_3 = 2d = 4 \end{cases}$, 解得 $a_1 = 3$, $d = 2$, 所以 $a_n = 2n+1$. $\displaystyle\sum_{i=2^{n-1}}^{2^n-1} a_i$, 项数为 $2^n-1 - 2^{n-1} + 1 = 2^{n-1}$, 首项 $a_{2^{n-1}} = 2\cdot2^{n-1}+1 = 2^n+1$, 末项 $a_{2^n-1} = 2(2^n-1)+1 = 2^{n+1}-1$, 所以和为 $2^{n-1} \cdot \dfrac{(2^n+1)+(2^{n+1}-1)}{2} = 2^{n-1} \cdot \dfrac{3\cdot2^n}{2} = 3 \cdot 2^{2n-1}$.
(2) (Ⅰ) 由题意, 当 $2^{k-1} \le n \le 2^k-1$ 时, $b_k < a_n$. 取 $n = 2^{k-1}$, 则 $b_k < a_{2^{k-1}} = 2^k+1$. 当 $2^{k-2} \le n \le 2^{k-1}-1$ 时, $a_n < b_k$. 取 $n = 2^{k-1}-1$, 则 $a_{2^{k-1}-1} = 2(2^{k-1}-1)+1 = 2^k-1 < b_k$. 所以 $2^k-1 < b_k < 2^k+1$.
(Ⅱ) 由 (Ⅰ) 知 $1 < b_1 < 3$, $3 < b_2 < 5$, $7 < b_3 < 9$, $15 < b_4 < 17$, 猜测 $b_n = 2^n$. 若公比 $q > 2$, 则 $b_n = b_1 q^{n-1} > 2^{n-1}$, 但 $2^{n-1} > 2^n-1$ 不恒成立, 矛盾. 若公比 $q < 2$, 则 $b_n = b_1 q^{n-1} < 3 \cdot 2^{n-1}$, 但 $3\cdot2^{n-1} < 2^n+1$ 不恒成立, 矛盾. 故 $q=2$, $b_1 = 2$, 所以 $b_n = 2^n$. 前 $n$ 项和 $S_n = \dfrac{2(1-2^n)}{1-2} = 2^{n+1} - 2$.
\end{solution}
\end{question}

\begin{question}[points=16]
已知函数 $f(x) = \left(\dfrac{1}{x} + \dfrac{1}{2}\right)\ln(x+1)$.

(1) 求曲线 $y = f(x)$ 在 $x = 2$ 处切线的斜率;

(2) 当 $x > 0$ 时, 证明: $f(x) > 1$;

(3) 证明: $\dfrac{5}{6} < \ln(n!) - \left(n+\dfrac{1}{2}\right)\ln n + n \le 1$.
\begin{solution}
\textbf{答案：}(1) $\dfrac{1}{3} - \dfrac{\ln 3}{4}$; (2) 证明见解析; (3) 证明见解析

\textbf{详解：}(1) $f(x) = \dfrac{\ln(x+1)}{x} + \dfrac{\ln(x+1)}{2}$, $f'(x) = \dfrac{1}{x(x+1)} - \dfrac{\ln(x+1)}{x^2} + \dfrac{1}{2(x+1)}$, 所以 $f'(2) = \dfrac{1}{6} - \dfrac{\ln3}{4} + \dfrac{1}{6} = \dfrac{1}{3} - \dfrac{\ln3}{4}$.
(2) 要证 $x>0$ 时 $f(x) > 1$, 即 $\ln(x+1) > \dfrac{2x}{x+2}$. 令 $g(x) = \ln(x+1) - \dfrac{2x}{x+2}$, $g'(x) = \dfrac{1}{x+1} - \dfrac{4}{(x+2)^2} = \dfrac{x^2}{(x+1)(x+2)^2} > 0$, 所以 $g(x)$ 在 $(0,+\infty)$ 上递增, $g(x) > g(0) = 0$, 故 $\ln(x+1) > \dfrac{2x}{x+2}$, 即 $f(x) > 1$.
(3) 设 $h(n) = \ln(n!) - \left(n+\frac{1}{2}\right)\ln n + n$, $n \in \mathbb{N}^*$. 则 $h(n+1)-h(n) = 1 - \left(n+\frac{1}{2}\right)\ln\left(1+\frac{1}{n}\right)$. 由 (2) 知 $\left(n+\frac{1}{2}\right)\ln\left(1+\frac{1}{n}\right) > 1$, 所以 $h(n+1)-h(n) < 0$, $\{h(n)\}$ 递减, 故 $h(n) \le h(1) = 1$. 下证 $h(n) > \dfrac{5}{6}$. 令 $\varphi(x) = \ln x - \dfrac{(x+5)(x-1)}{4x+2}$, $x>0$, 则 $\varphi'(x) = \dfrac{(x-1)^2(1-x)}{x(2x+1)^2}$, 当 $0<x<1$ 时 $\varphi'(x)>0$, 当 $x>1$ 时 $\varphi'(x)<0$, 所以 $\varphi(x) \le \varphi(1)=0$, 即 $\ln x \le \dfrac{(x+5)(x-1)}{4x+2}$ 恒成立. 则 $h(n)-h(n+1) = \left(n+\frac{1}{2}\right)\ln\left(1+\frac{1}{n}\right) - 1 \le \left(n+\frac{1}{2}\right) \cdot \dfrac{\left(\frac{6+\frac{1}{n}}{\frac{2}{n}}\right)}{\frac{4}{n}+2} = \dfrac{1}{12}\left(\frac{1}{n-1} - \frac{1}{n}\right)$. 累加得 $h(2)-h(n) < \frac{1}{12}\left(1-\frac{1}{n}\right)$, 而 $h(2) = 2-\ln2$, 所以 $h(1)-h(n) < \ln2 - 1 + \frac{3}{2} + \frac{1}{12}\left(1-\frac{1}{n}\right) < \ln2 + \frac{1}{2} + \frac{1}{12} = \ln2 + \frac{7}{12} < \frac{5}{6}$, 故 $h(n) > \frac{5}{6}$. 综上, $\frac{5}{6} < h(n) \le 1$.
\end{solution}
\end{question}

